\vspace{2.5pt plus 3pt minus 0.5pt}\noindent\hspace*{1.2em}\kbvnum{(6)}\hspace{0.55em}“ನಕ್ಷತ್ರ ಸಮುದಾಯದಲ್ಲಿ ಧ್ರುವ ನಕ್ಷತ್ರವು ಹೇಗೋ, ಹೊಸ ಗನ್ನಡಭಾಷೆಯ ಕವಿ ಸಮುದಾಯದಲ್ಲಿ ದ. ರಾ. ಬೇಂದ್ರೆಯವರು ಹಾಗೆ.”\par

\noindent\hfill{\fontsize{9.5pt}{12.5pt}\selectfont\color{kbvgrey}— ಈ ವಾಕ್ಯದಲ್ಲಿರುವ ಉಪಮಾನ, ಉಪಮೇಯಗಳನ್ನು ತಿಳಿಸಿರಿ.}\par\vspace{3pt plus 2pt}

\vspace{2.5pt plus 3pt minus 0.5pt}\noindent\hspace*{1.2em}\kbvnum{(7)}\hspace{0.55em}ಇವುಗಳ ಸೂತ್ರ ಬರೆದು, ಉದಾಹರಣೆ ಸಹಿತ ವಿವರಿಸಿ.\par

\vspace{2.5pt plus 3pt minus 0.5pt}\noindent\hspace*{1.2em}\kbvnum{(ಅ)}\hspace{0.55em}ಉತ್ಪ್ರೇಕ್ಷಾಲಂಕಾರ, (ಬ) ಶ್ಲೇಷಾಲಂಕಾರ,\par

\vspace{2.5pt plus 3pt minus 0.5pt}\noindent\hspace*{1.2em}\kbvnum{(ಕ)}\hspace{0.55em}ವೃತ್ತಾನುಪ್ರಾಸ, (ಡ) ಛೇಕಾನುಪ್ರಾಸ.\par

\vspace{2.5pt plus 3pt minus 0.5pt}\noindent\hspace*{1.2em}\kbvnum{(8)}\hspace{0.55em}ಇವುಗಳಲ್ಲಿ ಕಂಡುಬರುವ ಭಾಷಾಲಂಕಾರವನ್ನು ತಿಳಿಸಿರಿ.\par

\vspace{2.5pt plus 3pt minus 0.5pt}\noindent\hspace*{1.2em}\kbvnum{(೧)}\hspace{0.55em}ತಾಯಿಯೇ ದೇವರು.\par

\vspace{2.5pt plus 3pt minus 0.5pt}\noindent\hspace*{1.2em}\kbvnum{(೨)}\hspace{0.55em}ವಿಶ್ವಕನ್ನಡ ಸಮ್ಮೇಳನದ ಅವಧಿಯಲ್ಲಿ ಮೈಸೂರು ನಗರವು ಇಂದ್ರಲೋಕದಂತೆ ಕಂಗೊಳಿಸುತ್ತಿತ್ತು.\par

\vspace{2.5pt plus 3pt minus 0.5pt}\noindent\hspace*{1.2em}\kbvnum{(೩)}\hspace{0.55em}ಜಾಲಿಯ ಮರವು ನೆರಳಲ್ಲ-ಹೆಣ್ಣೀಗೆ\par

\begin{kbvverse}
\noindent ತಾಯೀಯ ಮನೆಯು ಸ್ಥಿರವಲ್ಲ ॥
\end{kbvverse}

\vspace{2.5pt plus 3pt minus 0.5pt}\noindent\hspace*{1.2em}\kbvnum{(೪)}\hspace{0.55em}ಬಡತನದಲ್ಲಿ ಹುಟ್ಟಿ ಬೆಳೆದ ಮೋಕ್ಷಗುಂಡಂ ವಿಶ್ವೇಶ್ವರಯ್ಯ ನವರು ತಮ್ಮ ಸ್ವಂತ ಪರಿಶ್ರಮದಿಂದ ವಿದ್ಯಾರ್ಜನೆ ಹಾಗೂ ಜ್ಞಾನಾರ್ಜನೆ ಮಾಡಿ ಮೈಸೂರಿನ ದಿವಾನರಾದರು ಮತ್ತು ಪ್ರತಿಭಾವಂತ ಇಂಜನಿಯರ್‌ ಎನಿಸಿಕೊಂಡು ಲೋಕ ವಿಖ್ಯಾತರಾದರು ಕೈಕೆಸರಾದರೆ ಬಾಯಿ ಮೊಸರು ಎಂಬ ಮಾತು ಸುಳ್ಳಲ್ಲ.\par

\vspace{2.5pt plus 3pt minus 0.5pt}\noindent\hspace*{1.2em}\kbvnum{(೫)}\hspace{0.55em}ಹಾರುತಿದೇ ನೋಡಿರಲ್ಲಿ ನಮ್ಮ ಬಾವುಟ\par

\begin{kbvverse}
\noindent ಮೂರು ಬಣ್ಣದಿಂದ ನಲಿವ ನಮ್ಮ ಬಾವುಟ ।
\end{kbvverse}

\indent ಗಗನವೇರಿ ಭರದಿ ಕುಣಿವ ದಿವ್ಯ ಬಾವುಟ\par

\begin{kbvverse}
\noindent ಗಾಂಧಿ ತಾತ ಗಳಿಸಿಕೊಟ್ಟ ಭವ್ಯ ಬಾವುಟ ॥
\end{kbvverse}

\vspace{2.5pt plus 3pt minus 0.5pt}\noindent\hspace*{1.2em}\kbvnum{(೬)}\hspace{0.55em}ಮಹಾಪುರುಷನು ವಿಪತ್ತಿಗೆ ಒಳಗಾದಾಗ ಹಿಂದೆಂದೂ ಕಾಣದ ಉದಾತ್ತ ಗುಣವನ್ನು ಪ್ರಕಟಿಸುತ್ತಾರೆ; ಧೂಪವು ಬೆಂಕಿಯ ನಡುವೆ ಇದ್ದಾಗ ಲೋಕೋತ್ತರವಾದ ಪರಿಮಳವನ್ನು ಹರಡುತ್ತದೆ. -ರಸಗಂಗಾಧರ.\par

\vspace{2.5pt plus 3pt minus 0.5pt}\noindent\hspace*{1.2em}\kbvnum{(೭)}\hspace{0.55em}ಭೀಮಸೇನನು ಬಾಲ್ಯದಲ್ಲಿ ನದಿ ತೀರದಲ್ಲಿ ಆಟವಾಡುತ್ತಿದ್ದಾಗ ಆನೆಗಳನ್ನು ಕೈಯಿಂದ ಎತ್ತಿ ನೀರಿಗೆ ಎಸೆಯುತ್ತಿದ್ದನಂತೆ!\par

\kbvsubtitle{0}
